\newcommand{\DoNotLoadEpstopdf}{}
\documentclass[11pt,reqno]{amsart}
\usepackage[T1]{fontenc}
\usepackage{lmodern,amsmath,amssymb,amsthm,mathtools,mathrsfs,microtype}
\usepackage[margin=1.1in]{geometry}
\usepackage[hidelinks]{hyperref}
\hypersetup{pdftitle={Collision records, arithmetic, and raw local inversion in a triangular Lorentz gas},pdfauthor={Qian Qi}}
\setlength{\emergencystretch}{2em}
\allowdisplaybreaks[2]
\numberwithin{equation}{section}
\newtheorem{theorem}{Theorem}[section]
\newtheorem{lemma}[theorem]{Lemma}
\newtheorem{proposition}[theorem]{Proposition}
\newtheorem{corollary}[theorem]{Corollary}
\theoremstyle{definition}\newtheorem{definition}[theorem]{Definition}
\newtheorem{inputstatement}[theorem]{Input statement}
\theoremstyle{remark}\newtheorem{remark}[theorem]{Remark}
\newcommand{\R}{\mathbb R}\newcommand{\Z}{\mathbb Z}
\newcommand{\one}{\mathbf1}\newcommand{\dd}{\,\mathrm d}
\newcommand{\esssup}{\operatorname*{ess\,sup}}
\title[Collision records and raw local inversion]{Collision records, arithmetic, and raw local inversion in a triangular Lorentz gas}
\author{Qian Qi}
\date{October 15, 2026. A2-DYN, clearance-flux author revision}
\subjclass[2020]{37D50, 37A50, 60F05, 28A75}
\keywords{Lorentz gas, exact collision records, optical Jacobi matrices, paired endpoint traces, positive source variation, global roof--angle uniqueness, marked pressure, incidence-load concentration, unvisited near misses, collision flux}
\begin{document}
\raggedbottom
\begin{abstract}
We retain the exact-label raw pointwise local limit and local inversion
problem for actual returns in the original triangular Lorentz gas.
The optical, arithmetic and incidence-load proof chain is preserved.
For unvisited-disk near misses, we prove an exact one-copy flux
identity: their expected number per physical collision is $\delta/R$.
A uniform gap between the marked disk and the incident disks gives
a transverse clearance-strip estimate which remains valid at grazing
limits. With the retained homogeneous stable-curve growth input,
this yields an $O(\varepsilon^{1-a})$ bound, for $0\le a<1$, in an
explicit stable-test seminorm after arbitrarily many actual
collisions and over all guard depths. An exact marked-record
disintegration preserves the physical roof and every return and
displacement label. These results provide a geometric mechanism for
clearance-only source, not a replacement of the physical dynamics.
The transfer from these norms to collision-normalized target-roof
height, low-incidence-load shortage, and the complete raw pointwise
assertion are not claimed proved.
\end{abstract}
\maketitle

\section*{Introduction}
The new contributions are Theorems~\ref{thm:cf-flux},
\ref{thm:cf-transverse}, \ref{thm:cf-transfer}, and
\ref{thm:cf-record}. They address actual unvisited-disk clearance,
including trajectories having no incidence witness. The first two
and the exact record identity have direct physical geometric proofs;
the arbitrary-time stable-test bound has the separately stated
growth input. The pointwise bridge is not inferred from either an
integrated flux or a weak transfer estimate. The preceding
incidence payments remain part of the same article and are not
counted again as new results.


The retained proved results are global roof--angle
uniqueness, target-roof marked pressure relative to the explicit
geometric inputs, and collision-normalized positive payments for
deep, numerous and high-load incidence sources. The complete
exact-label raw pointwise local limit assertion
\eqref{eq:v76-target} is still the research target: its equivalent
remaining estimate \eqref{eq:v112-remaining} is not proved here.
The retained load and quadratic-count theorems strengthen the positive
source estimates within that same problem, not a replacement problem.

The pointwise raw local limit problem concerns the density at the
original target roof, uniformly in collision count, radius and exact
return labels. Control of a fixed-window average does not by itself
control that density. For the original positive microscopic paired
source $\mathsf U_B^{\rm micro}$ and its full nongrazing comparison
pool $P_{r_B}^{\rm ng}$, the estimate left by the preceding argument is
\begin{equation}\label{eq:v110-intro-original}
 \lim_{B\to\infty}\limsup_{m\to\infty}\sup_{R,n,k}m^2
 \esssup_{t\in\mathcal C_{m,R,n,k}(M)}
 (\mathsf U_{B,m}^{\rm micro}(t)
                      -K_BP_{r_B,m}^{\rm ng}(t))_+=0.
\end{equation}
This is equivalent, within the retained scalar reduction, to the
original pointwise assertion \eqref{eq:v76-target}. It is not yet
proved here. The retained source-side contribution is a quantitative payment of
a further positive part of this same source and a physical description
of every coefficient which survives that payment.

The distinction between source and capacity is essential. On a
segment with an actually visited grazing collision at one end, the
original paired minimum is zero. Its endpoint maximum, and hence its
comparison capacity, need not be zero. The optical and trace arguments
retained in Sections~\ref{sec:v109-action}--\ref{sec:v109-shortage}
pay this entire visited-grazing capacity without deleting unvisited
tangencies or positive critical and section-boundary traces.
They start from the actual broken-path action, use a positive
continuant denominator, extend its reduced Hessian continuously to
visited grazing, and integrate the inverse gradient along a fixed
Gauss angle. All these arguments and their physical hypotheses are
retained in the present article.

The source-side argument uses the total variation of the
weighted density on an entire physical segment. If $c^\varepsilon$
is its endpoint minimum and $\mathcal V^\varepsilon$ its variation,
we split, before any roof pushforward,
\[
 c^\varepsilon
   =\min\{c^\varepsilon,\gamma_B\mathcal V^\varepsilon\}
       +(c^\varepsilon-\gamma_B\mathcal V^\varepsilon)_+,
 \qquad \gamma_B=\max\{2,\varepsilon(B)^{-1/32}\}.
\]
The first term can be positive, even when both endpoint traces are
positive. Theorem~\ref{thm:v110-payment} pays its complete physical
sum by the original uncut weighted volume, with no residence lower
bound, word count or comparison reach assumption. The resulting
ordered height is $O(B^{-95/384})$.

For the second term, Theorem~\ref{thm:v110-persistence} shows that the
density stays within relative distance $\gamma_B^{-1}$ of its
endpoint minimum and the original guard stays positive at every
point of the closed segment. In particular, a single point where
the original guard vanishes pays the entire original paired
coefficient. This assertion is independent of the name of either
endpoint mechanism. It does not assert that a single incidence or
clearance test witnesses the guard throughout the segment.

Theorem~\ref{thm:v110-residual} retains the full nongrazing comparison
pool and gives an exact three-term positive identity for the
original microscopic source. The remaining estimate is printed in
\eqref{eq:v110-remaining-limit}, in the original essential-height
topology and the original order of limits. Neither nearly constant
density on a short segment nor finite-count exhaustion is identified
with a proof of this remaining estimate.

Appendix~\ref{sec:v110-optical-audit} reconstructs the signed corner
cofactor and proves a simultaneous multi-mark expansion. It also
states the simultaneous one-sided endpoint lemma separately and
makes the uniform constant in the trace modulus explicit. The
multi-mark bound is per trajectory; it is not a new weighted-pressure
theorem. The complete preceding article and historical derivations
are preserved byte-for-byte in the original repository archive \cite{Q110}.
The physical pairing model, regular derivative estimate and spectral
average estimates remain source-pinned inputs rather than consequences
of finite algebraic diagnostics.

\subsection*{From momentum fibers to the original target roof}
A geometric pressure bound on endpoint momentum fibers does not by
itself control a target-roof fiber: the momentum magnitudes of its
words need not agree. Theorem~\ref{thm:v111-unique} removes this
obstruction to summation. Convexity of the Cartesian broken-path
length gives two support inequalities between any two physical
trajectories of one lifted itinerary. On a common endpoint-gradient
ray these inequalities order their lengths strictly by gradient
magnitude. Thus a word supplies at most one point at a specified
roof and Gauss angle, without requiring a connected word domain.

Theorem~\ref{thm:v111-roof-pressure} combines this global fact with
the retained uniform geometric pressure and optical stable-cocycle
comparison. For any set $S$ of distinct interior marks it gives
\[
 \left\|\frac{dF_\#(\prod_{i\in S}c_i^{-1}\one_\Pi\nu_R^*)}{dt}
             \right\|_\infty\le C_\star K_\star^{|S|}.
\]
All original exact labels remain on the left. On the upper side,
a marked full word is reconstructed from its unmarked blocks and
its finitely many deleted physical step symbols. This is a weighted
summation theorem; it does not insert an uncontrolled count of words
or homogeneous strips. Its constants have no Gaussian $m^{-2}$
factor, a distinction essential to its use.

Corollary~\ref{cor:v111-witness-moments} sums changing incidence
witnesses by elementary symmetric polynomials. The original geometric
depth thresholds have bounded sum, independently of $m$.
Theorem~\ref{thm:v111-witness-payment} consequently pays every
persistent-source trajectory with an incidence witness beyond an
explicit $O(\log m)$ end layer, with normalized height
$C\varepsilon(B)(m+1)^{-\eta}$. A factorial bound also pays states
with $O(\log m/\log\log m)$ or more incidence witnesses. These
collision-dependent classifications are substituted only into bounds
valid at every finite count; all spectral windows and outer scales
remain fixed in their original order.

This leaves the original pointwise target, not a new topic, at
\eqref{eq:v111-remaining}. The remaining shallow incidence and
clearance source is not declared small by an unscaled pressure
estimate. Sections~\ref{sec:v111-global}--\ref{sec:v111-witness}
contain the new proofs; all preceding source and capacity arguments
remain in the article, with the historical packet pinned in the original repository
under \texttt{archive/v111/archive/v110} \cite{Q110}.


\subsection*{Joint incidence load and the two collision-end layers}
A witness at depth $d_i$ already spends the amount
$\alpha d_i$ of logarithmic incidence load, where
$\alpha=|\log q_0|/4$. Its additional smallness spends
$\log(2\varepsilon q_0^{d_i/4}/c_i)$. These contributions are
nonnegative and can be summed before selecting a witness. The
resulting variable $\mathscr L_\varepsilon$ obeys
\[
 \left\|\frac{dF_\#(e^{\lambda\mathscr L_\varepsilon}
                               \one_\Pi\nu_R^*)}{dt}\right\|_\infty
 \le C_\star\exp\left\{\frac{4K_\star\varepsilon}
                                  {e^{\alpha(1-\lambda)}-1}\right\},
 \qquad 0<\lambda<1.
\]
Theorem~\ref{thm:v112-load} proves this by a positive product
expansion on the original physical source. Optimization gives the
tail $C_\star\exp[-u+2\sqrt{b_\varepsilon u}]$, where
$b_\varepsilon=4K_\star\varepsilon/\alpha$.
At $u=(2+\eta)\log(m+1)$ this pays the normalized height, including
states with only one sufficiently severe shallow witness.

The depth multiset has at most two entries at every level.
Lemma~\ref{lem:v112-depth-poly} preserves this fact in the
coefficient of $\prod_{d\ge1}(1+zr^d)^2$, instead of bounding that
coefficient by a factorial involving only the total threshold sum.
It proves the quadratic exponent
$r^{\lfloor(s+1)^2/4\rfloor}$. The resulting
Theorem~\ref{thm:v112-count} reduces the witness-count threshold to
order $\sqrt{\log m}$. Sharpness is established for this depth
polynomial, not claimed for an unproved physical lower bound.

Theorem~\ref{thm:v112-payment} incorporates the new payment as a
positive restriction of the original source. Its retained part has
bounded total incidence depth and bounded total logarithmic severity
at the displayed collision scale. Clearance witnesses, including
unvisited tangency, remain distinct and may occur at every depth.
The full preceding article and all of its historical material are
preserved in \texttt{archive/v111} \cite{Q111}. The current article
retains every preceding active proof, promotes the geometric input
conventions explicitly, and supplies the unbounded-weight exhaustion
and marked-step indexing used by the target-roof sum.

\section{The unchanged physical problem}
Put $R_-=9/20$ and $R_+=47/100$, and let $R\in I=[R_-,R_+]$. Circular scatterers have centers in
$\Lambda=\Z(1,0)+\Z(1/2,\sqrt3/2)$. Write $T_R$ for the regular
collision map and $(\alpha,p)\in(\R/2\pi\Z)\times(-1,1)$ for its
coordinates; $p$ is the outgoing tangential velocity and
$\dd\nu=\dd\alpha\,\dd p/(4\pi)$. Physical flights satisfy
$\tau_-\le\tau_R\le\tau_+$, where $\tau_-=3/50$ and $\tau_+$
is the fixed finite-horizon bound.

The section is specified here to fix every event in the argument.
Put $h_0=1/200$. Let $Y_R$ be the union of the four open coordinate
squares of half-side $h_0$ centered at
\[
 (0,0),\quad(\pi/6,-1/2),\quad
 (-\theta_R,\sin\theta_R),\quad
 (\pi/3-\theta_R,\sin\theta_R).
\]
Here $\theta_R=2\arctan q_R$, and $q_R$ is the unique root in
$(49/100,53/100)$ of
\[
 8Rq(1+q^2)-4q(1-q^2)+\sqrt3(1-6q^2+q^4)=0.
\]
Set
\[
 U=\{|\alpha-\pi/6|<3/50,\ |p|<3/20\},\qquad
 Y_R^*=Y_R\cup U,\qquad c_*=\frac{91}{10000\pi}.
\]
The rectangles are disjoint and $\nu(Y_R^*)=c_*$. We always use
$\nu_R^*=c_*^{-1}\nu|_{Y_R^*}$, not a renormalization of a smaller
source.

At $m$ flights, let $L=L_{m,R}$ be their total length,
$K^{\rm c}_{m,R}$ their lattice displacement, and
$A_{m,R}=\sum_{j=0}^{m-1}\one_{Y_R^*}\circ T_R^j$. The event is
\begin{equation}\label{eq:v76-event}
 \Pi_{n,k,m,R}=Y_R^*\cap\{K^{\rm c}_{m,R}=k,\ A_{m,R}=n\}
                         \cap T_R^{-m}Y_R^*.
\end{equation}
The initial visit is counted; the terminal visit is tested separately.
On this event the actual $n$th return has collision count $m$ and
roof $L$. Denote its density by $p_{n,R}(k,m,t)$.

Write $c_j=\sqrt{1-p_j^2}$ and let $\Delta_j$ be the least clearance
of flight $j$ from any third disk, excluding its two incident disks.
Keep the original smooth function $0\le\Theta\le1$, zero on
$(-\infty,1]$ and one on $[2,\infty)$. With $q_0=47/53$,
$d_j=\min(j,m-j)$ and $d'_j=\min(j,m-1-j)$, set
\begin{align}
 H^\varepsilon&=
 \prod_{j=0}^{m}\Theta\!\left(\frac{c_j}{\varepsilon q_0^{d_j/4}}\right)
 \prod_{j=0}^{m-1}\Theta\!\left(\frac{\Delta_j}{\varepsilon q_0^{d'_j/4}}\right),
 \qquad W^\varepsilon=1-H^\varepsilon,\label{eq:v76-guard}\\
 b^\varepsilon(t)\dd t
   &=L_\#(W^\varepsilon\one_{\Pi_{n,k,m,R}}\nu_R^*).
 \label{eq:v76-source}
\end{align}
In particular first incidence and first clearance outside any selected
critical neighborhoods are still present. No section-distance factor
has been inserted into this guard.

\subsection{The arithmetic reference and the target}
For $y=(k_1,k_2,t,n)$ put
\[
 \mu_R=(0,0,\bar\tau_R,c_*),\quad
 Z=(y-m\mu_R)/\sqrt m,\quad
 \mathcal C_{m,R,n,k}(M)=\{t:|Z|\le M\}.
\]
We retain the finite spectral-peak reference, including peaks close to
but not on the unit circle. To specify its meaning, the collision
record observable is $(\kappa_1,\kappa_2,\tau,\one_{Y_R^*})$.
Let $\lambda_j(R,\vartheta)$ and $\mathsf P_j(R,\vartheta)$ be its
continued simple twisted eigenvalue and projection in each of the
fixed peak charts, and let $\chi_j$ be their frequency partition.
At the real maximum $a_j$ of $\log|\lambda_j|$ put
\[
 z_j=\lambda_j(R,a_j),\quad
 \mu_j=i^{-1}\nabla\log\lambda_j(R,a_j),\quad
 H_j=-\nabla^2\log\lambda_j(R,a_j),
\]
with $\operatorname{Re}H_j\ge aI_4$ for a common $a>0$. If
$\eta_R=\one_{Y_R^*}$ and $\ell$ denotes mass evaluation, put
$B_j=\chi_j(R,a_j)\ell(M_{\eta_R}\mathsf P_j(R,a_j)(\eta_R\nu))$.
Define
\begin{align}
 G_H(z)&=(2\pi)^{-4}\int_{\R^4}e^{-iv\cdot z-v^{\mathsf T}Hv/2}\dd v,
 \notag\\
 \mathcal L_{m,R}(y)&=\operatorname{Re}\left[
 c_*^{-1}\sum_j e^{-ia_j\cdot y}z_j^m B_j
      G_{H_j}\left((y-m\mu_j)/\sqrt m\right)\right].
 \label{eq:v76-reference}
\end{align}
Only the charts active at $R$ are summed. This is the same transition
reference as in the retained spectral theorem. The existence and
uniform estimates for these scalar data belong to that analytic
input, not to the transport argument below. No frequency branch or
arithmetic zero class is deleted.

The governing assertion remains
\begin{equation}\label{eq:v76-target}
 \sup_{R,\,1\le n\le m,\,k\in\Z^2}
 \esssup_{t\in\mathcal C_{m,R,n,k}(M)}
 |m^2p_{n,R}(k,m,t)-\mathcal L_{m,R}(k_1,k_2,t,n)|\longrightarrow0.
\end{equation}
Let $\mathcal E_M$ denote this error's collision limsup. An empty
central target contributes zero. The order of the uniformity is
$\sup_R\esssup_t$, so exceptional density sets may depend on $R$.



\section{The physical source and the inherited analytic inputs}
\label{sec:v109-inputs}
We specify the inputs to the new geometric argument, separating the
existence of the original physical source from the endpoint regularity
that will be proved below. The complete historical proofs and their
qualifications are retained in the accompanying archive.

For fixed $m,R,n,k$, let $\Theta_{\rm phys}$ be the original disjoint
incident-component pairing-sheet space. It is the countable Borel
sum of the maximal regular components of the original exact event,
parametrized by
\[
 (u,\theta)=(F,\arg\nabla_qF),\qquad q=(\alpha_0,\alpha_m),
 \qquad a(\theta)<u<b(\theta),\quad \ell=b-a>0.
\]
The sheet identity, and not just its angular coordinate, is included
in $\theta$. Integration $d\theta$ means the original one-copy
physical sum with angular Lebesgue measure. The regular components,
their Borel endpoint maps, their physical multiplicity convention,
and their common coarea versions are the retained pairing-space
input. This article does not replace them by symbolic collision words.
Every component carries the same exact event \eqref{eq:v76-event}.
Put
\begin{equation}\label{eq:v109-density}
 \begin{split}
 g&=\nabla_qF=R(-p_0,p_m),\quad z=|g|,
 \quad I_m=1+\sum_{i=1}^{m-1}q_0^{d_i}/c_i,\\
 \omega&=|\partial_{p_0}\alpha_m|^{-1},\quad
 \chi=|\det D_q(F,\arg g)|,\quad
 A=\frac{\omega}{4\pi c_*\chi},\quad f^\varepsilon=AW^\varepsilon.
 \end{split}
\end{equation}
The density formula is used initially on the regular interior.
Once its endpoint traces have been established, the original pairing
is
\begin{equation}\label{eq:v109-paired}
 f_-^\varepsilon=f^\varepsilon(a+,\theta),\quad
 f_+^\varepsilon=f^\varepsilon(b-,\theta),\quad
 c^\varepsilon=\min(f_-^\varepsilon,f_+^\varepsilon),\quad
 d\pi^{\varepsilon,+}=c^\varepsilon\one_{\{a<u<b\}}d\theta\,du.
\end{equation}
Write $p_\vee=\max(f_-^\varepsilon,f_+^\varepsilon)$ and
$e_\partial=|f_-^\varepsilon-f_+^\varepsilon|$.
Thus $p_\vee=c^\varepsilon+e_\partial$.

\begin{inputstatement}[Regular variation and completed physical averages]
\label{input:v109-analytic}
The inherited regular-point optical calculation gives, almost
 everywhere on each regular segment,
\begin{equation}\label{eq:v109-regular-variation}
 |\partial_u f^\varepsilon|\leq
 C(1+\varepsilon^{-1})AI_m/z.
\end{equation}
Let $\mathsf T_\Delta^\varepsilon$ be the complete paired roof
density restricted to $0<\ell<\Delta$, with no additional selector,
and let
\[
 v^0(t)\,dt=F_\#(I_m z^{-1}\one_\Pi\nu_R^*),\qquad
 A_rf(t)=(2r)^{-1}\int_{t-r}^{t+r}f(u)\,du.
\]
For every previously fixed $\varepsilon,r,\Delta,\rho>0$,
\begin{align}
 \limsup_m\sup_{R,n,k}m^2\|A_r\mathsf T_\Delta^\varepsilon\|_\infty
 &\leq C\{\varepsilon^{1/16}+(1+\varepsilon^{-1})\Delta\},
 \label{eq:v109-paired-average}\\
 \limsup_m\sup_{R,n,k}m^2\|A_\rho v^0\|_\infty&\leq C.
 \label{eq:v109-volume-average}
\end{align}
The constants are independent of the previously fixed widths.
The coarea identity identifies the full physical sum of $AI_m/z$
with $v^0$. Fixed-count local finiteness and common roof versions
are retained. These are averaged, not pointwise, bounds.
\end{inputstatement}
These are the positive pairing, uncut volume and weighted optical
variation statements in Section 2 of the frozen article \cite{Q108}.
Their exact historical source labels are recorded in the proof ledger.
The new proof
uses \eqref{eq:v109-regular-variation} only on regular interiors and
supplies its integrable passage to the physical endpoints.

\subsection{The original microscopic source}
The outer schedule is
\begin{equation}\label{eq:v109-schedule}
 \varepsilon(B)=A_0B^{-1/12},\qquad r_0=c_0^{\rm out}\varepsilon(B)^4,
 \qquad\delta(B)=\tfrac34r_0,\qquad\Delta_B=8\delta(B).
\end{equation}
The fixed constants $A_0,c_0^{\rm out}$ are the inherited constants;
$c_0^{\rm out}$ is written this way to avoid confusion with the
initial incidence $c_0$.
Keep the original remainder selector $H_B$, the paid critical-arm
family $\mathscr C_B$, and the adaptive paid family $\mathscr D_B$
of the historical microscopic reduction. Their memberships are
Borel, disjointized by first physical membership, and their total
outer costs are at most $C\varepsilon(B)^{-1/32}$.
For a previously fixed $0<h_B<\delta(B)$ put
\begin{equation}\label{eq:v109-original-micro}
 \begin{split}
 \widetilde H_B&=H_B\one_{(\mathscr C_B\cup\mathscr D_B)^c}
                       \one_{\{0<\ell<h_B\}},\\
 \mathsf U_B^{\rm micro}(t)
   &=\int c^{\varepsilon(B)}\one_{\{a<t<b\}}
                           \widetilde H_B(\Psi(\theta,t))\,d\theta.
 \end{split}
\end{equation}
No sign cancellation, deletion of an endpoint class, or change of
section occurs in this definition. Here $\Psi$ denotes the original
physical parametrization.

\begin{inputstatement}[The scalar microscopic reduction]
\label{input:v109-scalar}
For every fixed $M$, define
\[
 \mathcal H_M(T)=\limsup_m\sup_{R,n,k}m^2
          \esssup_{t\in\mathcal C_{m,R,n,k}(M)}T(t).
\]
The historical exact positive reduction gives
\begin{equation}\label{eq:v109-scalar-input}
 \mathcal E_M\leq\mathcal H_M(\mathsf U_B^{\rm micro})
                              +C_MB^{-1/384}.
\end{equation}
The assertion \eqref{eq:v76-target} is equivalent to
$\mathcal H_M(\mathsf U_B^{\rm micro})\to0$ for every fixed $M$.
All widths and costs are fixed at the outer stage before the
complete physical sum, the essential suprema, and the collision
limsup; $B$ is outermost.
\end{inputstatement}
This is the unchanged theorem in historical module 255. The
fixed-positive-residence estimate in module 253 and the earlier
paid-arm estimates are retained as its inputs. A bound uniform in
a previously fixed positive residence is not used at a
collision-dependent residence.

\input{core/260_optical_jacobi_and_endpoint_trace_completion}
\input{core/261_grazing_capacity_payment_and_nongrazing_shortage}
\input{core/262_positive_variation_payment}
\input{core/263_persistent_guard_and_nongrazing_shortage}
\input{core/265_global_roof_angle_uniqueness}
\input{core/266_target_roof_marked_pressure}
\input{core/267_uniform_incidence_witness_summation}
\input{core/268_joint_incidence_load_concentration}
\input{core/269_positive_low_load_source_reduction}
\appendix
\section{An exact flux identity for unvisited near misses}
\label{sec:cf-flux}
The incidence load does not see a trajectory whose original guard is
witnessed only by clearance. We introduce a counting observable on the
\emph{same} physical flights, not a new collision map. Throughout this
section let
\[
 g_*=1-2R_+=3/50,\qquad \delta_*=g_*/4=3/200,
 \qquad 0<\delta\le\delta_*.
\]
The symbol $\delta_*$ here is a one-flight geometric constant, not the
outer microscopic width $\delta(B)$. All disks still have radius $R$.

For a regular outgoing collision $x$, write its flight as
$q(x)+tv(x)$, $0\le t\le\tau(x)$. Let $\mathcal Z_\delta(x)$ be the
set of lattice centers other than the two incident centers for which
\[
 0<\min_{0\le t\le\tau(x)}|q(x)+tv(x)-z|-R<\delta,
 \qquad N_\delta(x)=|\mathcal Z_\delta(x)|.
\]
This counts near misses, not collisions. In particular neither the
collision index, the physical roof, nor the section-visit count is
incremented by $N_\delta$. A center is counted once on a flight.
Tangencies may be assigned value zero on their collision-area null set;
no assertion about a positive pairing-endpoint trace follows from this
convention.

\begin{lemma}[An interior foot and uniform endpoint separation]
\label{lem:cf-foot}
Every center counted by $N_\delta$ has its perpendicular foot in the
open physical flight. If $t_z$ is the time of that foot, then
\begin{equation}\label{eq:cf-gap}
 t_z\ge g_*-\delta,\qquad \tau-t_z\ge g_*-\delta.
\end{equation}
The candidate centers, after translating the initial incident center
to zero, belong to the finite set
\[
 \{z\in\Lambda\setminus\{0\}:|z|\le\tau_++2R_++\delta_*\}.
\]
Write $N_*$ for its cardinality. This constant is independent of the
itinerary length and of all exact record labels.
\end{lemma}
\begin{proof}
At either incident point, the distance to a third center is at least
$1-R$, since distinct lattice centers have distance at least one.
Its distance to the third disk is therefore at least $1-2R\ge g_*$.
A minimum below $\delta\le g_*/4$ cannot occur at an endpoint. The
nearest point on a line segment is then its perpendicular foot, whose
distance from $z$ is at most $R+\delta$. By the reverse triangle
inequality its distance along the line from either endpoint is at
least $(1-R)-(R+\delta)\ge g_*-\delta$. Finally a foot is at distance
at most $R+\tau_+$ from the initial center. Adding its distance from
$z$ gives the finite-set bound. This proof uses the actual segment,
including both of its endpoints, rather than an infinite-line
substitute for the clearance event.
\end{proof}

Put $v_\vartheta=(\cos\vartheta,\sin\vartheta)$ and
$v_\vartheta^\perp=(-\sin\vartheta,\cos\vartheta)$. For a mark $z$
define the signed line coordinate
$b=(z-q)\cdot v_\vartheta^\perp$. Translate the marked center to zero.
For $\sigma\in\{-1,1\}$ and $0<u<\delta_*$ the oriented line with
$b=\sigma(R+u)$ has closest point
$-\sigma(R+u)v_\vartheta^\perp$. Follow this line backwards to its
last actual collision. Reduce that collision by a lattice translation
to the initial fundamental disk, and denote the resulting outgoing
state and marked-center displacement by
\[
 X_R(\vartheta,\sigma,u),\qquad Z_R(\vartheta,\sigma,u).
\]
These maps use straight free flight until an \emph{actual} disk is
met. There is no reflection at the closest point. The finite-horizon
assumption ensures actual collisions in both directions, apart from
null singular parameters on which these maps can be defined
arbitrarily. The interval from the preceding collision to the following
one remains the original physical flight.

\begin{theorem}[One-copy clearance flux and exact multiplicity]
\label{thm:cf-flux}
For every nonnegative Borel function $a$ of an outgoing state and a
marked-center displacement,
\begin{equation}\label{eq:cf-flux}
 \int\sum_{z\in\mathcal Z_\delta(x)}a(x,z)\,d\nu(x)
 =\frac{1}{4\pi R}\sum_{\sigma=\pm1}
   \int_0^{2\pi}\int_0^\delta
       a(X_R(\vartheta,\sigma,u),Z_R(\vartheta,\sigma,u))
                                  \,du\,d\vartheta.
\end{equation}
In particular,
\begin{equation}\label{eq:cf-exact-mean}
                         \int N_\delta\,d\nu=\delta/R.
\end{equation}
Both statements hold for every $R\in[R_-,R_+]$. The multiplicity on
the left is the number of physically distinct unvisited disks, not
an upper bound for the number of collision words.
\end{theorem}
\begin{proof}
We first verify the local measure, including its radius factor. Write
$q=R(\cos\alpha,\sin\alpha)$ and
$\vartheta=\alpha+\arcsin p$. At fixed $\vartheta$,
\[
 \left.\frac{\partial b}{\partial\alpha}\right|_\vartheta
     =-R\cos(\vartheta-\alpha),\qquad
 \frac{\partial\vartheta}{\partial p}=(1-p^2)^{-1/2}.
\]
Consequently
\begin{equation}\label{eq:cf-jacobian}
 \left|\det D_{(\alpha,p)}(\vartheta,b)\right|=R,
 \qquad d\nu=\frac{d\vartheta\,db}{4\pi R}.
\end{equation}
The two signs $b=\sigma(R+u)$ both have $|db|=du$.

There is also a global multiplicity assertion to prove. A marked
physical flight maps to its unique oriented line relative to its
marked center. Conversely, the closest point of every line just
described is in the free region: its distance to any other lattice
center is at least $1-R-u>R$. In fact the disk of radius $R+\delta_*$
around the marked center does not intersect any other actual disk,
since $2R+\delta_*<1$. The central disk is missed because the line
has impact magnitude greater than $R$. Following the line in both
directions therefore gives a physical flight whose incident disks
are different from the marked disk. Its foot is interior and its
clearance from that disk is $u$. This is inverse to the marked-flight
map, modulo the prescribed lattice translation. Multiple near misses
on one flight give different marked copies; they are neither lost nor
identified.

For measure-theoretic precision, partition the line parameters by the
actual backward incident disk and the forward incident disk, and
exclude their tangent boundaries. These partitions are Borel; on each
regular piece the disk-intersection formulas and their inverses are
smooth. Tangency boundaries satisfy the line equation $|b_z|=R$.
For any fixed alternative center their equation in $(\vartheta,u)$
is a graph in $u$ (with one of the two fixed signs), hence has planar
Lebesgue measure zero. There are countably many centers. The
corresponding singular collision sets are also null by
\eqref{eq:cf-jacobian} on regular charts. Exhaust these charts by
relatively compact subcharts and disjointize them. Applying the local
Jacobian formula to bounded functions on finite unions, then using
monotone convergence for both chart unions and $a\wedge n$, proves
\eqref{eq:cf-flux} without a prior integrability assumption. Setting
$a=1$ gives $2(2\pi)\delta/(4\pi R)=\delta/R$.
\end{proof}

\begin{corollary}[All physical depths with the original normalization]
\label{cor:cf-depth-mass}
For arbitrary $0<\delta_j\le\delta_*$, $0\le j<m$,
\[
 \int\sum_{j=0}^{m-1}N_{\delta_j}(T_R^jx)\,d\nu(x)
       =\frac1R\sum_{j=0}^{m-1}\delta_j.
\]
For every original exact event $\Pi$ the same integral with
$\one_\Pi\,d\nu_R^*$ is at most
$(Rc_*)^{-1}\sum_j\delta_j$. In particular, if $2\varepsilon\le\delta_*$
and $r=q_0^{1/4}$, then
\begin{equation}\label{eq:cf-all-depth-mass}
 \nu_R^*\!\left(\Pi\cap
   \{\exists j:\Delta_j<2\varepsilon r^{d'_j}\}\right)
       \le\frac{4\varepsilon}{R_-c_*(1-r)}.
\end{equation}
\end{corollary}
\begin{proof}
The collision measure $\nu$ is invariant under the original elastic
collision map. This also follows from the oriented-line Jacobian
\eqref{eq:cf-jacobian}, whose magnitude is the same at both incident
circles, and specular reflection, which preserves boundary flux.
Thus apply \eqref{eq:cf-exact-mean} at each $j$. The original initial
restriction has density $\one_{Y_R^*}/c_*\le1/c_*$ with respect to
$\nu$; further exact labels only decrease the positive integral.
Outside the null exact-tangency set, existence of a clearance witness
implies $N_{2\varepsilon r^{d'_j}}\ge1$. There are at most two flight
indices of each nonnegative depth, so
$\sum_j2\varepsilon r^{d'_j}\le4\varepsilon/(1-r)$. No independence
or return-conditioned invariant measure is introduced.
\end{proof}
The estimate \eqref{eq:cf-all-depth-mass} is an integrated mass
estimate. Multiplication of it by $m^2$ at fixed $\varepsilon(B)$
is not a proof about a target-roof density.

\section{Clearance strips transverse to physical test curves}
\label{sec:cf-transverse}
A measure bound alone gives no short-curve control. The positive
separation between an unvisited disk and the incident disks supplies
a transverse derivative which remains nonzero even as an incident
angle approaches grazing. We specify the time orientation explicitly:
forward near misses are tested on increasing-slope curves, and
\emph{past-flight} near misses on decreasing-slope curves.

Use physical arclength $s=R\alpha$ increasing counterclockwise and
$p=\sin\varphi$. Fix a finite $U>0$ and put
\begin{equation}\label{eq:cf-curve-constants}
 \begin{split}
 \kappa_-&=1/R_+,\qquad \kappa_+=1/R_-,\\
 D_0&=\kappa_++U+(1-R_+)^{-1},\\
 \ell_0&=\pi/(2D_0),\qquad
 a_0=\kappa_-(g_*-\delta_*),\\
 C_{\rm geo}&=2N_*\sqrt{1+U^2}/a_0.
 \end{split}
\end{equation}
All lengths below are Euclidean arclength in $(s,\varphi)$, not
configuration-space arclength. A test is a connected graph over an
$s$ interval, of length at most $\ell_0$. Let $E_\delta^+$ be the
forward physical near-miss event $\{N_\delta\ge1\}$, and let
$E_\delta^-$ be the analogous event on the preceding actual flight.

\begin{theorem}[A grazing-uniform, one-flight strip bound]
\label{thm:cf-transverse}
For every $C^1$ test graph $W^+$ with
$0\le d\varphi/ds\le U$, and every $C^1$ test graph $W^-$ with
$-U\le d\varphi/ds\le0$,
\begin{equation}\label{eq:cf-strip}
 |W^+\cap E_\delta^+|\le C_{\rm geo}\delta,
 \qquad |W^-\cap E_\delta^-|\le C_{\rm geo}\delta.
\end{equation}
The bound is uniform in $R$, in the homogeneity strip of the test,
and in its distance to the grazing boundary. Arbitrary Borel
restrictions, including exact record restrictions, can be added to
the sets on the left. Regular pieces are understood with their
physical multiplicity one.
\end{theorem}
\begin{proof}
On $W^+$ the direction angle is $\vartheta=s/R+\varphi(s)$.
For a fixed candidate center $z\ne0$ set
\[
 t_z=(z-q)\cdot v_\vartheta,\qquad
 b_z=(z-q)\cdot v_\vartheta^\perp.
\]
Differentiation in the increasing $s$ coordinate gives the signed
identity
\begin{equation}\label{eq:cf-transverse-derivative}
 b'_z=-\cos\varphi-t_z\vartheta',\qquad
 \kappa_-\le\vartheta'\le\kappa_++U.
\end{equation}
At a physical near miss, Lemma~\ref{lem:cf-foot} yields
$t_z\ge g_*-\delta_*$. Hence $-b'_z\ge a_0$ there. It is the
$t_z\vartheta'$ term, not a lower bound on $\cos\varphi$, which
prevents degeneration.

We must not multiply this estimate by an uncontrolled number of
components. Choose a continuous argument $\beta(s)=\arg(z-q(s))$
on the graph. Since $|z-q|\ge1-R_+$, one has
$|\beta'|\le(1-R_+)^{-1}$. Moreover
\[
 (\beta-\vartheta)'\le (1-R_+)^{-1}-\kappa_-<0,
\]
because $R_+<1/2$. Its total variation is at most
$D_0|W^+|\le\pi/2$. Thus the set $t_z>0$, equivalently
$\cos(\beta-\vartheta)>0$, is at most one interval: two distinct
positive-cosine intervals are separated by a negative interval of
angular length $\pi$. On the interval $t_z>0$, formula
\eqref{eq:cf-transverse-derivative} says that $b_z$ strictly decreases.
Each of the two impact bands $(R,R+\delta)$ and
$(-R-\delta,-R)$ therefore has multiplicity at most one in this
coordinate. Restricting to the portions which are actual near misses
can break them into pieces, but cannot increase that multiplicity.
By the one-dimensional change of variables on this monotone map,
\[
 \int_{\{s:\ z\text{ is a physical near miss}\}}ds
                 \le 2\delta/a_0.
\]
Since $d\ell\le\sqrt{1+U^2}\,ds$, summing over the finite candidate
set from Lemma~\ref{lem:cf-foot} proves the first bound. This sum
counts only centers in one physical flight. It does not count
iterated collision words or homogeneous continuity components.

For the second bound, the collision time-reversal involution is
$\mathcal I(s,\varphi)=(s,-\varphi)$. The outgoing velocity at
$\mathcal I x$ reverses the incoming velocity at $x$. Consequently
its next straight flight is the reversal of the previous actual
flight at $x$. This proves
\begin{equation}\label{eq:cf-reversal}
 E_\delta^-=\mathcal I E_\delta^+=T_RE_\delta^+
\end{equation}
on the regular collision space. The involution is an isometry in
$(s,\varphi)$ and sends $W^-$ to an increasing-slope graph.
Apply the first bound. Borel restrictions decrease lengths. Taking
monotone limits of truncated regular test pieces gives the same
bound at grazing limits; an exactly tangent parameter contributes
zero transverse length by the strict derivative above. This last
statement does not delete any paired-source boundary trace in the
older coarea construction.
\end{proof}

\begin{corollary}[Bounded multipliers on stable tests]
\label{cor:cf-multiplier}
If $|h|\le H$ pointwise on regular collision states, then for every
$W^-$ as above and every bounded Borel test $\psi$,
\[
 \left|\int_{W^-}\one_{E_\delta^-}h\psi\,d\ell\right|
       \le H\|\psi\|_\infty
                   \min\{|W^-|,C_{\rm geo}\delta\}.
\]
This is a bound for bounded functions with specified representatives.
It is not an assertion that multiplication by the indicator is
bounded on an entire space of anisotropic distributions.
\end{corollary}
\begin{proof}
Bound the integrand by $H\|\psi\|_\infty\one_{E_\delta^-}$ and
use Theorem~\ref{thm:cf-transverse} together with the trivial length
bound. Representatives on singular sets are fixed to zero.
\end{proof}

\section{Clearance gains after arbitrarily many actual collisions}
\label{sec:cf-transfer}
We next give a consequence which is uniform in the full collision
count, rather than only a one-flight estimate. Its additional input
is the retained stable-curve growth bound, stated explicitly here.
It is not the stronger target-roof local limit estimate.

\begin{inputstatement}[Retained homogeneous curve growth]
\label{input:cf-growth}
Use a family $\mathcal W_R^s$ of short homogeneous admissible stable
graphs in the above physical coordinates, with slopes in $[-U,0]$,
length at most $\ell_0$, and the inherited common curvature bound.
For every $W\in\mathcal W_R^s$ and $l\ge0$, its regular inverse
image has a countable partition $\mathcal G_l(W)$ into such graphs,
up to arclength-null sets, with one-copy inverse branches and
\begin{equation}\label{eq:cf-growth}
 \sum_{V\in\mathcal G_l(W)}\sup_V J_VT_R^l\le C_g.
\end{equation}
Here $J_VT_R^l$ is the Euclidean arclength derivative along $V$;
$C_g$ is independent of $R,l,W$. At $l=0$ the partition is $\{W\}$.
\end{inputstatement}
This is the curve-growth interface of historical module 241, not its
whole-cell pressure conclusion. Length thresholds can be decreased:
subdividing a pre-existing short component into lengths at most
$\ell_0$ multiplies its Jacobian upper sum by at most a fixed number,
independent of $l$. The precise external interface is discussed in
Appendix~\ref{sec:cf-audit}. We retain this input separately from the
unconditional one-flight geometry of
Theorems~\ref{thm:cf-flux} and \ref{thm:cf-transverse}.

For fixed representatives of bounded Borel functions put
\[
 \mathcal S_a(f)=\sup_{W\in\mathcal W_R^s}
       |W|^{-a}\int_W|f|\,d\ell,\qquad 0\le a<1.
\]
This absolute stable-test seminorm is specified here to avoid
identifying it with a full anisotropic distribution norm. With respect
to the invariant measure $\nu$, the ordinary density transfer is
$\mathcal L_Rh=h\circ T_R^{-1}$. For the original record observable
$\mathbf g_R=(\kappa_1,\kappa_2,\tau,\one_{Y_R^*})$ and real
$\xi=(\xi_1,\xi_2,\xi_t,\xi_n)$, define
\[
 \mathcal L_{R,\xi}h(x)=
 e^{i\xi\cdot\mathbf g_R(T_R^{-1}x)}h(T_R^{-1}x).
\]
The three discrete frequencies may be taken modulo $2\pi$; the
physical-roof frequency is real and unbounded. This definition does
not assert boundedness of this operator on an unstated Banach space.
On bounded functions its modulus is simply $|h|\circ T_R^{-1}$.

\begin{theorem}[A uniform future-propagated clearance gain]
\label{thm:cf-transfer}
Under Input~\ref{input:cf-growth}, let $|h|\le H$ pointwise on the
regular collision space, with $h=0$ off $E_\delta^-$. Then for every
$l\ge0$, real $\xi$, and $0\le a<1$,
\begin{align}
 \int_W|\mathcal L_{R,\xi}^{\,l}h|\,d\ell
  &\le H\min\{|W|,C_g C_{\rm geo}\delta\},
                                  \label{eq:cf-future-length}\\
 \mathcal S_a(\mathcal L_{R,\xi}^{\,l}h)
  &\le H(C_g C_{\rm geo}\delta)^{1-a}.
                                  \label{eq:cf-future-test}
\end{align}
The constants do not depend on the entire future collision word or
on any of the four real frequencies. The bounds also hold for any
bounded multiplier of modulus at most one applied at arbitrary
intermediate collisions, with the same support at the marked time.
\end{theorem}
\begin{proof}
On regular branches positivity and $|e^{i\xi\cdot\mathbf g_R}|=1$
give
$|\mathcal L_{R,\xi}^{\,l}h(x)|
 \le H\one_{E_\delta^-}(T_R^{-l}x)$.
Change variables separately on the inverse branches of $W$. The
full nonnegative sum, completed by monotone convergence, satisfies
\[
 \int_W|\mathcal L_{R,\xi}^{\,l}h|\,d\ell
 \le H\sum_{V\in\mathcal G_l(W)}
       (\sup_V J_VT_R^l)\,|V\cap E_\delta^-|
 \le HC_{\rm geo}\delta C_g.
\]
The pointwise bound by $H$ supplies the other term in the minimum.
For $x>0$ and $D>0$, splitting into $x\le D$ and $x\ge D$ gives
$x^{-a}\min(x,D)\le D^{1-a}$. Apply this with
$x=|W|$ and $D=C_gC_{\rm geo}\delta$ to obtain
\eqref{eq:cf-future-test}. The modulus of a product of additional
multipliers is at most one, so precisely the same positive upper
bound applies. No stochastic independence, wordwise supremum in
place of a physical integral, or Fourier inversion has been used.
\end{proof}

\begin{corollary}[The complete depth schedule in stable tests]
\label{cor:cf-all-depth-tests}
Let $f_0=\one_{Y_R^*}/c_*$. For $0\le j<m$ set
\[
 h_j=\one_{E_{\delta_j}^-}\mathcal L_{R,\xi}^{\,j+1}f_0,
 \qquad
 f_{m,j}=\mathcal L_{R,\xi}^{\,m-j-1}h_j.
\]
Then
\begin{equation}\label{eq:cf-mark-transfer}
 \mathcal S_a(f_{m,j})\le
  c_*^{-1}(C_gC_{\rm geo}\delta_j)^{1-a}.
\end{equation}
If $2\varepsilon\le\delta_*$ and
$\delta_j=2\varepsilon r^{d'_j}$, the terminal collision density
restricted to the union of all original clearance witnesses has
absolute stable-test seminorm at most
\begin{equation}\label{eq:cf-scheduled-tests}
 \frac{2(C_gC_{\rm geo})^{1-a}(2\varepsilon)^{1-a}}
             {c_*(1-r^{1-a})},\qquad 0\le a<1.
\end{equation}
This holds uniformly over $m,R$ and real record frequencies. Further
original exact labels and terminal-section restrictions may be
retained on the left by positive domination.
\end{corollary}
\begin{proof}
The initial density and all of its twisted iterates have modulus at
most $1/c_*$. By \eqref{eq:cf-reversal}, the past-flight event at
collision $j+1$ is exactly the forward clearance mark on flight $j$.
Thus the exponents $j+1$ and $m-j-1$ insert the mark at the actual
collision time, including $j=0$ and $j=m-1$ with zero remaining
future steps. Theorem~\ref{thm:cf-transfer} proves
\eqref{eq:cf-mark-transfer}. The union indicator is bounded by the
sum of its physical one-flight indicators. Subadditivity of
$\mathcal S_a$ and at most two indices at each depth give
\[
 \sum_{j=0}^{m-1}\delta_j^{1-a}
      \le\frac{2(2\varepsilon)^{1-a}}{1-r^{1-a}}.
\]
This proves \eqref{eq:cf-scheduled-tests}. Arbitrary exact-record
selectors are bounded by one on the same trajectories, so they
cannot increase the positive majorant. They are not replaced by
independent weights, and no probability measure is renormalized.
\end{proof}

The estimate is genuinely uniform in the full number of physical
collisions and in the full scheduled set of clearance marks. Its
norm is nevertheless a stable-test norm on terminal collision space,
not the essential height of a roof density. In particular, the absence
of frequency growth in this absolute estimate is not decay as
$|\xi_t|\to\infty$ and does not justify integration over that
unbounded frequency. Such decay, or another rigorously justified
pointwise bridge, is still required for the original raw target.

\section{Exact-record disintegration at an unvisited mark}
\label{sec:cf-record}
The flux coordinates also give a direct representation of the
remaining roof problem which does not discard the exact return,
displacement or collision information. We state it as an equality
of measures first, because a conditional angular roof kernel need
not have a density without proof.

For $0\le j<m$ define the initial state
\[
 X_{R,j}(\vartheta,\sigma,u)=T_R^{-j}X_R(\vartheta,\sigma,u)
\]
on its regular set, with a fixed null-set convention otherwise.
The complete actual record from this state is used. Displacements
are computed in the original lattice basis, relative to its initial
lift; translating the marked disk to zero changes none of these
relative displacements. For almost every $u$, define the positive
Borel measure on the original roof line
\begin{equation}\label{eq:cf-kernel}
 \mathcal K_{m,R,n,k,j,u}(A)=\sum_{\sigma=\pm1}
  \int_0^{2\pi}\one_{\Pi_{n,k,m,R}}(X_{R,j})
             \one_{\{L_{m,R}(X_{R,j})\in A\}}\,d\vartheta.
\end{equation}
It has mass at most $4\pi$. No density or angular derivative bound
is included in this definition.

\begin{theorem}[Disintegration of a physically marked exact record]
\label{thm:cf-record}
For every Borel set $A\subset\R$,
\begin{equation}\label{eq:cf-record-disintegration}
 \int\one_\Pi N_\delta(T_R^jx)
                 \one_{\{L_{m,R}\in A\}}\,d\nu_R^*
   =\frac{1}{4\pi R c_*}\int_0^\delta
                       \mathcal K_{m,R,n,k,j,u}(A)\,du.
\end{equation}
The identity also holds against every nonnegative Borel roof test.
All singular boundary sets in these integrals use the original
collision-area null-set convention, not deletion of paired boundary
coefficients.
\end{theorem}
\begin{proof}
Replace $d\nu_R^*$ by $c_*^{-1}d\nu$ inside the integral, since
$\Pi\subset Y_R^*$. Change the initial collision coordinate to
$y=T_R^jx$ using invariance of $\nu$. Apply
Theorem~\ref{thm:cf-flux} with
\[
 a(y,z)=\one_\Pi(T_R^{-j}y)
               \one_{\{L_{m,R}(T_R^{-j}y)\in A\}}.
\]
This is a bounded nonnegative function on the original physical
states, independent of which near-miss mark is selected once $y$
is fixed. The resulting integral is exactly
\eqref{eq:cf-record-disintegration}. Countable regular partitions and
Tonelli's theorem justify the integration over $u$; their exceptional
parameters have two-dimensional measure zero. For arbitrary
nonnegative roof tests use monotone simple-function approximation.
In particular the factor is $1/(4\pi R c_*)$, with no second
section normalization and no new conditioning on the near miss.
\end{proof}

\subsection{A sufficient pointwise bridge, with its status explicit}
Fix $M$. The following measure domination is one concrete, stronger
sufficient condition for paying all original clearance witnesses:
there exists $C_M$ such that, for all sufficiently large $m$, all
$R,n,k,j$ and almost every $0<u<\delta_*$,
\begin{equation}\label{eq:cf-bridge-input}
 \mathcal K_{m,R,n,k,j,u}\big|_{\mathcal C_{m,R,n,k}(M)}
       \le C_M m^{-2}\,dt.
\end{equation}
An almost-everywhere set may depend on the fixed $R,m,n,k,j$;
$C_M$ must be common. This is an \emph{unproved sufficient condition},
not a newly assumed premise of the organizing raw local limit
assertion. It asserts absence of singular conditional roof mass as
well as an upper density bound. Neither Theorem~\ref{thm:cf-flux}
nor Theorem~\ref{thm:cf-transfer} establishes it.

\begin{proposition}[The exact normalization a pointwise bridge must pay]
\label{prop:cf-bridge}
If \eqref{eq:cf-bridge-input} holds, the original roof measure
restricted to any clearance witness, with $2\varepsilon\le\delta_*$,
is dominated on the central target by a density $d_{\rm clr}$ with
\begin{equation}\label{eq:cf-bridge-consequence}
 \sup_{R,n,k}m^2\|d_{\rm clr}\|_{L^\infty(\mathcal C_{m,R,n,k}(M))}
       \le\frac{C_M\varepsilon}{\pi R_-c_*(1-r)}.
\end{equation}
The same upper bound holds for any positive source already dominated
by that original collision measure, restricted to these witnesses.
\end{proposition}
\begin{proof}
The witness-union indicator is at most
$\sum_jN_{2\varepsilon r^{d'_j}}(T_R^jx)$ off the fixed null set.
Apply \eqref{eq:cf-record-disintegration} to every term, then the
measure domination \eqref{eq:cf-bridge-input}. The resulting density
bound is $(4\pi Rc_*)^{-1}C_Mm^{-2}\sum_j2\varepsilon r^{d'_j}$.
The sum is at most $4\varepsilon/(1-r)$, proving the stated constant.
Positive domination preserves the bound for the indicated source.
At the original outer schedule this would tend to zero after the
collision limsup and then $B\to\infty$; no window has been changed.
\end{proof}
On a regular physical branch, the new parametrization has the exact
chain rule
\[
 D X_{R,j}=D T_R^{-j}(X_R)\,D X_R,\qquad
 D(L_{m,R}\circ X_{R,j})=D L_{m,R}(X_{R,j})\,D X_{R,j}.
\]
The one-flight lower bound for $|b'_z|$ does not lower-bound this
long-time roof derivative. Inverse-iterate growth, cancellations and
critical angular roofs must be controlled separately. A bound after
integration in $u$ can allow angular folds which a uniform bound on
each $\mathcal K_{j,u}$ would exclude. No supremum is passed through
that disintegration here.

A proof of the stronger condition \eqref{eq:cf-bridge-input} is not
necessary if a suitably summed kernel bound or comparison-capacity
argument can be proved instead. The proposition only makes explicit
which normalization and which measure must be controlled. It is not
counted as an unconditional closure or a second payment of a source
already paid in revision 112.

\subsection{Compatibility with the complete original source}
All identities through \eqref{eq:v112-positive-identity} and the full
pool $P_{r_B}^{\rm ng}$ are unchanged. We have not subtracted the
clearance source from $\mathsf U_B^{\rm low}$ using only a weak norm.
Consequently the exact remaining target is still
\eqref{eq:v112-remaining}, with the original order of limits.
The new results treat the mechanism which is invisible when
$J_\varepsilon=\varnothing$: genuine unvisited-disk near misses.
They do not pay low-severity shallow incidence without clearance,
prove abundance of the full nongrazing comparison pool, or settle
positive critical-entrance and section-cut pairing traces.

For precision concerning a minor referee point, the inequality
$\alpha f(N)<u_m$ already proved in
Corollary~\ref{cor:v112-support} implies
$N<2\sqrt{u_m/\alpha}$ on the retained source. The independent
quadratic count theorem is retained as a finite-count theorem, not
presented as another missing payment for that already restricted
source. The term ``low'' in all inherited source symbols means low
\emph{incidence} load only. In the monotone exhaustion of
Proposition~\ref{prop:v112-exhaustion}, every marked weight is
explicitly assigned zero on the omitted collision-area null set;
its extended integral is taken only over the stated regular charts.
This convention never sets a positive one-sided pairing trace to zero.

\section{Scope of the retained comparison statements}
\label{sec:v109-scope}
For clarity, we record the precise scope of the earlier allocation
statements. None of their algebraic maximality assertions is used
as a substitute for the physical estimate in
Theorem~\ref{thm:v109-residual}.

\begin{proposition}[Fixed-window scalar maximality and old-family dominance]
\label{prop:v109-maximal}
Fix a Borel comparison subspace, $0<r<\Delta$, and $K>0$.
Consider nonnegative measurable weights $w_+,w_-$ with
$w_++w_-\leq1$, supported on the respective one-sided window
conditions in \eqref{eq:v109-pool}. Among scalar accepted heights
$Z(t)$ satisfying
\[
 0\leq Z(t)\leq U(t),\qquad
 Z(t)\leq K\int(w_++w_-)p_\vee\,d\theta,
\]
the greatest is $V_{r,K}(t)=\min\{U(t),KP_r(t)\}$.
The allocation is global at the roof and does not couple individual
source sheets to comparison sheets.

For a compatible old fixed family with window width $r_{\rm old}$,
physical overlap $N$ and block cost $\kappa$, the explicit conditions
\[
 0<r\leq r_{\rm old},\qquad \ell_{\rm comparison}<\Delta,
 \qquad K=N\kappa
\]
imply $V^{\rm old}(t)\leq V_{r,N\kappa}(t)$ on the same original
source and comparison space.
\end{proposition}
\begin{proof}
The support and one-copy conditions give $w_++w_-\leq Q_r$.
Hence every $Z$ is bounded by the stated minimum. On $P_r>0$ choose
$w_\pm=V_{r,K}q_\pm/(KP_r)$; use zero weights otherwise.
This is a measurable attaining certificate. The separate source
multiplier is $\beta_{r,K}$ defined in
Section~\ref{sec:v109-shortage}.
For the old family, retain only active source blocks at $t$.
Their comparison supports lie in the $r$-pool by the displayed
width condition. Summing the blockwise budget and using physical
overlap at most $N$ gives $V^{\rm old}\leq N\kappa P_r$;
disjoint original source blocks also give $V^{\rm old}\leq U$.
The minimum proves the assertion, with countable families obtained
by monotone convergence on their common roof versions.
\end{proof}

\begin{proposition}[Strictness relative to roof-independent whole-interval blocks]
\label{prop:v109-strictness}
In the interval-source model, the roofwise certificate class is
strictly larger than the old class of roof-independent angular
source blocks that requires a comparison interval to contain the
translate of the \emph{entire} source interval by $[0,r]$ or
$[-r,0]$. This statement compares certificate classes and does not
assert a Lorentz realization.
\end{proposition}
\begin{proof}
Let $s=1/1000$ and $r=1/10$. On two disjoint angular sets of measure
one give both endpoint traces value one. Let their physical intervals
be $(0,s)$ and $(3s/8,r+7s/8)$, and select only the first as source.
Take $\Delta>r+s/2$ and $K=1$. For $3s/8<t<7s/8$ the second
interval contains $[t,t+r]$, so $P_r(t)\geq U(t)=1$.
Neither physical interval has length $r+s$, which is necessary for
the old translated-whole-source inclusion. Angular subdivision does
not alter these intervals. Thus the old class cannot allocate a
nonzero whole-source block while the roofwise class accepts the
source on the indicated roof interval.
\end{proof}

\subsection{Limit order and the remaining estimate}
The constructions in this article are evaluated first at finite
$m,R,n,k$ with $B,r_B,h_B,K_B$ already fixed. Positive physical sums
are completed before taking the essential roof supremum, the radius
and exact-label suprema, and the collision limsup. Only then is $B$
sent to infinity. The assertion that every positive-length segment
contains a sufficiently small window at a fixed count supplies no
uniform modulus for this order.

The remaining estimate is precisely
\[
 \lim_{B\to\infty}\limsup_{m\to\infty}\sup_{R,n,k}m^2
 \esssup_{t\in\mathcal C_{m,R,n,k}(M)}
 \bigl(\mathsf U_{B,m}^{\rm micro}(t)
                 -K_BP_{r_B,m}^{\rm ng}(t)\bigr)_+=0.
\]
No inference to this estimate is made from bounded wordwise optical
weights, local integrability, a finite diagnostic, or a positive
finite-count incidence. Theorem~\ref{thm:v109-residual} proves the
entire visited-grazing payment in this order; all remaining positive
classes in the displayed expression are still part of the original
pointwise problem.

\input{core/264_optical_cofactors_and_simultaneous_endpoints}
\section{Dependency, reference and normalization audit}
\label{sec:cf-audit}
The direct clearance proofs use the original positive disk separation,
the uniform finite-horizon hypothesis, circular curvature, elastic
reflection and elementary change of variables. The arbitrary-time
stable-test estimate additionally uses Input~\ref{input:cf-growth}.
Neither uses uniform pressure, a twisted spectral gap, or the
unproved localized-shortage estimate to prove its own premise.

\subsection{The precise imported curve estimate}
In the inspected 10 October 2013 author version of
Demers--Zhang~\cite{DZUniform}, Lemma 3.2(b), printed page 14,
bounds the sum of suprema of the forward arclength Jacobians on
homogeneous inverse-image pieces, uniformly for $W$ and for maps
in the stipulated class. The geometric-family statement is Theorem
2.7 in that version; the varying-arclength qualification is Remark
2.9(b). The inherited module 241's numbering must therefore be
matched to a specified version rather than silently treated as
verified. Its stronger finite-testing and pressure argument is not
needed for Theorem~\ref{thm:cf-transfer}.

For this paper the coordinates, Euclidean arclength Jacobian,
short-curve curvature bound, homogeneous partition, fixed subdivision
length, and map-family constants must all agree with
Input~\ref{input:cf-growth}. Changing boundary perimeter with $R$
requires the coordinate transport and its uniform equivalence;
changing the subdivision threshold changes $C_g$ by a bounded
factor. These interfaces remain explicit retained hypotheses here,
rather than a claim of an independent specialist reconstruction.

In the inspected 19 September 2022 author version of
Baladi--Demers~\cite{BDPressure}, Lemma 2.2, printed pages 7--8,
is a distortion statement across stable curves in the same
homogeneous backward cell. Its closure convention excludes true
singularity curves. It is not a four-record, roof-conditioned
$m^{-2}$ estimate. The present proofs make no such inference.

\subsection{Checks which enter the proofs}
The radius factor in \eqref{eq:cf-jacobian} is $R$, not $R^2$ or
$1/R$. The two impact signs and the full direction circle give the
factor $4\pi\delta$ before collision normalization. Exactly one
section factor $c_*^{-1}$ occurs in
\eqref{eq:cf-record-disintegration}. The foot has physical time
coordinate $t_z$, so the sign in
\eqref{eq:cf-transverse-derivative} is a minus for both the boundary
motion and the direction motion. Increasing-slope curves, not
forward stable curves, are used for the forward event.

The uniform lower bounds used above are strictly positive on the
whole closed radius interval:
\[
 g_*-\delta_*=9/200,\qquad
 \kappa_--(1-R_+)^{-1}=600/2491,\qquad
 a_0=9/94.
\]
The occurrence of $N_*$ in the curve estimate, but not in the exact
flux identity, is deliberate: the former is a one-flight union
bound on a test curve, whereas the latter is a global one-copy
change of section. Empty future blocks are included by $l=0$.
The denominator $1-r^{1-a}$ is finite only for $a<1$; no estimate
at the endpoint $a=1$ is claimed.

The source identities in the retained proof bodies are unchanged.
All positive sums are completed before essential suprema and the
collision limsup. Taking $B\to\infty$ in a stable-test estimate
would not by itself permit taking that limit in the target-roof
norm. Fixed-$m$ simulations and symbolic identities check algebra
and implementation only; they do not certify either the inherited
uniform dynamical inputs or the remaining raw pointwise theorem.

\subsection{What is retained, and what has not been reconstructed}
The complete active revision-112 manuscript and its ten active core
proof bodies are retained in this manuscript, and an unmodified
copy of the downloaded active packet is supplied in
\texttt{provenance/v112-active}. The original repository additionally
contains the recursive revision-111 archive, with source tree
\texttt{790e627efe3a89a4549123d7e9acacf24089c287}.
That recursive archive was not exported by the completed CI artifact
available for this revision. Its identity is preserved; it has not
been represented as a locally reconstructed, independently audited
set of 267 historical proofs. The included source manifest and
preservation audit distinguish byte-verified local preservation from
remote tree pinning. In particular the historical spectral,
arithmetic, induced-return and four-coordinate inverse-problem
inputs are not replaced by the new clearance lemmas.

\begin{thebibliography}{9}
\bibitem{Q111} Q. Qi,
\emph{Collision records, arithmetic, and raw local inversion in a
triangular Lorentz gas}, A2-DYN revision 111, October 14, 2026.
Complete source packet preserved in \texttt{archive/v111}.
Author commit \texttt{8c617cafa5e5493e26817401e866b2415522f7c0}.

\bibitem{Q110} Q. Qi,
\emph{Collision records, arithmetic, and raw local inversion in a
triangular Lorentz gas}, A2-DYN revision 110, October 14, 2026.
Complete source packet preserved in \texttt{archive/v111/archive/v110}.
Author commit \texttt{0abce3fcfbc735ff5cc6d085279abea5f32a7a22}.
\bibitem{Q109} Q. Qi,
\emph{Collision records, arithmetic, and raw local inversion in a
triangular Lorentz gas}, A2-DYN revision 109, October 14, 2026.
Complete source packet preserved in \texttt{archive/v111/archive/v110/archive/v109}.
Author commit \texttt{ab0ad79b9c101b82c80e33511529588d101a45c0}.
\bibitem{Q108} Q. Qi,
\emph{Collision records, arithmetic, and raw local inversion in a
triangular Lorentz gas}, A2-DYN revision 108, October 14, 2026.
Frozen manuscript and complete historical supplement, preserved
byte-for-byte in \texttt{archive/v111/archive/v110/archive/v109/archive/v108}. Author commit
\texttt{2144665f89362ca80116f1999d99c80875833dcf}.
\bibitem{DZUniform} M. F. Demers and H.-K. Zhang,
\emph{A functional analytic approach to perturbations of the Lorentz gas},
Communications in Mathematical Physics \textbf{324} (2013), 767--830.
The inspected author version is dated 10 October 2013.
\bibitem{BDPressure} V. Baladi and M. F. Demers,
\emph{Thermodynamic formalism for dispersing billiards},
Journal of Modern Dynamics \textbf{18} (2022), 415--493.
The inspected author version is dated 19 September 2022.
\end{thebibliography}
\end{document}
